
le nombre des polynômes \(\chi(F(n))\) est au plus égal au nombre de composantes connexes de \(T\) (EGA III 7.9.5) et c'est un petit lemme que si pour un \(t\in T\), \(F_t\) est localement libre de rang \(p\), alors \(F\) l'est aussi au-dessus d'un voisinage ouvert de \(t\) dans \(T\), (voir EGA IV 11.3.10, où on fait \(X=Y\), \(f=\mathrm{id}\)). Or, il résulte des hypothèses et de (EGA III 7.7.6) que quitte à couper \(T\) encore plus fin, on peut trouver une suite \(\sigma\) de sections \(\sigma_1,\ldots,\sigma_r\in H^0({O}_X(1))\) qui est \(F\)-régulière. Donc (ii) résulte de la propriété de semi-continuité supérieure de la fonction \(t\longmapsto h^0(F_{t,i}(-1))\) où pour \(t\in T\) et \(0\le i\le r\), \(F_{t,i}\) désigne la restriction de \(F_t\) au schéma des zéros de \(\sigma_1,\ldots,\sigma_i\) (EGA III 7.7.5).

L'implication (ii) \(\Rightarrow\) (iii) résulte immédiatement de 1.11. L'implication (iii) \(\Rightarrow\) (iv) résulte de 1.3 (iii) en prenant \(M=\max(\chi(F_K(m))\) et \(E={O}_X(-m)^{\oplus M}\).

Supposons que \(\mathcal F\) satisfasse à (iv); il existe des suites exactes
\[
0\longrightarrow F'_K\longrightarrow E_K\longrightarrow F_K\longrightarrow0
\]
et l'ensemble des polynômes \(\chi(F_K(n))\) est fini. Par hypothèse la famille des classes des \(E_K\) est limitée; donc par l'implication (i) \(\Rightarrow\) (ii) déjà prouvée, l'ensemble des polynômes \(\chi(E_K(n))\) est fini et il existe une suite d'entiers \((b)\) telle que tous les \(E_K\) (avec \(K\) algébriquement clos) soient des \((b)\)-faisceaux. Par suite, l'ensemble des polynômes \(\chi(F'_K(n))\) est fini et en vertu de 1.6, tous les \(F'_K\) sont des \((b)\)-faisceaux. Donc en appliquant à la famille des classes des \(F'_K\) l'implication (ii) \(\Rightarrow\) (iv) déjà prouvée, on trouve (v).

Supposons enfin que \(\mathcal F\) satisfasse à (v), prouvons que \(\mathcal F\) est limitée. On peut supposer \(E\) (resp. \(E'\)) de la forme \({O}_X(-m)^{\oplus M}\). En effet,
